% GENERATED FILE. Edit canonical lessons, then run npm run generate:books.
% canonical-source-hash: fnv1a64:e074bd88455c1beb
% canonical-lessons: IT-C399-account, IT-C399-password, IT-C399-utente, IT-C399-accedere, IT-C399-registrarsi

\chapter{Il mio account}
\label{ch:it-il-mio-account}
\hlchaptermodality{399}

\begin{chapteropening}
\textbf{By the end of this chapter:} \emph{I can manage an online account: the user, the password, signing up and logging in.}

Complete the last lesson of chapter 399: I can manage an online account: the user, the password, signing up and logging in.
\end{chapteropening}

\section[l'account]{l'account — the online account}

\label{lesson:IT-C399-account}



\begin{quote}
Before the new one: say the Italian for the loss, then the Italian for to go bankrupt.
\end{quote}

\subsection*{You'll want to know: l'account}
\textbf{l'account} — \textquotedblleft{}the online account\textquotedblright{}.

\begin{quote}
Se non hai un account, devi crearne uno.
\end{quote}

\begin{quote}
\emph{If you do not have an account, you have to create one.}
\end{quote}

Masculine, invariable: \emph{gli account}. A bank account is \emph{il conto}.

\begin{cousinweb}
\emph{account} was borrowed from English \emph{account} (EtymDB), which is usually traced to Old French \emph{acont}, built on \emph{conter} \textquotedblleft{}to count\textquotedblright{}, from Latin \emph{computāre} \textquotedblleft{}to reckon\textquotedblright{}.
\end{cousinweb}

\begin{grammarlens}[title={What you've built}]
One more word for telling, explaining and describing at B1.
\end{grammarlens}

\subsection*{Your turn}
\begin{itemize}
\raggedright
  \item \emph{Say it:} \emph{l'account}
  \item \emph{Say it:} \emph{l'account}, once more
  \item \emph{Say it:} a sentence of your own with \emph{l'account}
  \item \emph{Recall:} say the Italian for the loss, then the Italian for to go bankrupt, then say \emph{l'account} again
\end{itemize}

\begin{tcolorbox}[breakable,colback=teal!4,colframe=teal!35!black,title={Before you move on}]
What does l'account mean? (\textquotedblleft{}The online account\textquotedblright{}.) Say it once more.
\end{tcolorbox}
\section[la password]{la password — the password}

\label{lesson:IT-C399-password}



\begin{quote}
Before the new one: say the Italian for to go bankrupt, then the Italian for the online account.
\end{quote}

\subsection*{You'll want to know: la password}
\textbf{la password} — \textquotedblleft{}the password\textquotedblright{}.

\begin{quote}
Se dimentichi la password, clicca qui.
\end{quote}

\begin{quote}
\emph{If you forget your password, click here.}
\end{quote}

Feminine in Italian; less often \emph{la parola d'ordine}.

\begin{cousinweb}
\emph{password} was borrowed from English \emph{password} (EtymDB), built from \emph{pass} and \emph{word}: the word that lets you pass a guard.
\end{cousinweb}

\begin{grammarlens}[title={What you've built}]
One more word for telling, explaining and describing at B1.
\end{grammarlens}

\subsection*{Your turn}
\begin{itemize}
\raggedright
  \item \emph{Say it:} \emph{la password}
  \item \emph{Say it:} \emph{la password}, once more
  \item \emph{Say it:} a sentence of your own with \emph{la password}
  \item \emph{Recall:} say the Italian for to go bankrupt, then the Italian for the online account, then say \emph{la password} again
\end{itemize}

\begin{tcolorbox}[breakable,colback=teal!4,colframe=teal!35!black,title={Before you move on}]
What does la password mean? (\textquotedblleft{}The password\textquotedblright{}.) Say it once more.
\end{tcolorbox}
\section[l'utente]{l'utente — the user}

\label{lesson:IT-C399-utente}



\begin{quote}
Before the new one: say the Italian for the online account, then the Italian for the password.
\end{quote}

\subsection*{You'll want to know: l'utente}
\textbf{l'utente} — \textquotedblleft{}the user\textquotedblright{}.

\begin{quote}
Il sito ha più di un milione di utenti.
\end{quote}

\begin{quote}
\emph{The website has more than a million users.}
\end{quote}

Also a customer of a service: \emph{gli utenti del treno}. One form for men and women.

\begin{cousinweb}
\emph{utente} is a learned borrowing from Latin \emph{ūtēns, ūtentis} \textquotedblleft{}using\textquotedblright{} (Lewis \& Short), present participle of \emph{ūtī} \textquotedblleft{}to use\textquotedblright{}.
\end{cousinweb}

\begin{grammarlens}[title={What you've built}]
One more word for telling, explaining and describing at B1.
\end{grammarlens}

\subsection*{Your turn}
\begin{itemize}
\raggedright
  \item \emph{Say it:} \emph{l'utente}
  \item \emph{Say it:} \emph{l'utente}, once more
  \item \emph{Say it:} a sentence of your own with \emph{l'utente}
  \item \emph{Recall:} say the Italian for the online account, then the Italian for the password, then say \emph{l'utente} again
\end{itemize}

\begin{tcolorbox}[breakable,colback=teal!4,colframe=teal!35!black,title={Before you move on}]
What does l'utente mean? (\textquotedblleft{}The user\textquotedblright{}.) Say it once more.
\end{tcolorbox}
\section[accedere]{accedere — to log in, to access; to enter}

\label{lesson:IT-C399-accedere}



\begin{quote}
Before the new one: say the Italian for the password, then the Italian for the user.
\end{quote}

\subsection*{You'll want to know: accedere}
\textbf{accedere} — \textquotedblleft{}to log in, to access; to enter\textquotedblright{}.

\begin{quote}
Per accedere, inserisci il tuo nome e la password.
\end{quote}

\begin{quote}
\emph{To log in, enter your name and your password.}
\end{quote}

\emph{Accedere a} — to get into: \emph{accedere al conto}, \emph{accedere all'università}.

\begin{cousinweb}
\emph{accedere} is a learned borrowing from Latin \emph{accēdere} \textquotedblleft{}to approach, to come near\textquotedblright{} (EtymDB), built from \emph{ad} \textquotedblleft{}to\textquotedblright{} and \emph{cēdere} \textquotedblleft{}to go\textquotedblright{}.
\end{cousinweb}

\begin{grammarlens}[title={What you've built}]
One more word for telling, explaining and describing at B1.
\end{grammarlens}

\subsection*{Your turn}
\begin{itemize}
\raggedright
  \item \emph{Say it:} \emph{accedere}
  \item \emph{Say it:} \emph{accedere}, once more
  \item \emph{Say it:} a sentence of your own with \emph{accedere}
  \item \emph{Recall:} say the Italian for the password, then the Italian for the user, then say \emph{accedere} again
\end{itemize}

\begin{tcolorbox}[breakable,colback=teal!4,colframe=teal!35!black,title={Before you move on}]
What does accedere mean? (\textquotedblleft{}To log in, to access; to enter\textquotedblright{}.) Say it once more.
\end{tcolorbox}
\section[registrarsi]{registrarsi — to sign up, to register}

\label{lesson:IT-C399-registrarsi}



\begin{quote}
Before the new one: say the Italian for the user, then the Italian for to log in.
\end{quote}

\subsection*{You'll want to know: registrarsi}
\textbf{registrarsi} — \textquotedblleft{}to sign up, to register\textquotedblright{}.

\begin{quote}
Se ti registri oggi, il primo mese è gratis.
\end{quote}

\begin{quote}
\emph{If you sign up today, the first month is free.}
\end{quote}

Without \emph{-si}, \emph{registrare} is to record: \emph{registrare un video}.

\begin{cousinweb}
\emph{registrarsi} is built on \emph{registro} \textquotedblleft{}register\textquotedblright{}, from Medieval Latin \emph{registrum}, from Late Latin \emph{regesta} \textquotedblleft{}a list\textquotedblright{}, from \emph{regerere} \textquotedblleft{}to carry back, to record\textquotedblright{} (Lewis \& Short), built on \emph{gerere} \textquotedblleft{}to carry\textquotedblright{}.
\end{cousinweb}

\begin{grammarlens}[title={What you've built}]
One more word for telling, explaining and describing at B1.
\end{grammarlens}

\subsection*{Your turn}
\begin{itemize}
\raggedright
  \item \emph{Say it:} \emph{registrarsi}
  \item \emph{Say it:} \emph{registrarsi}, once more
  \item \emph{Say it:} a sentence of your own with \emph{registrarsi}
  \item \emph{Recall:} say the Italian for the user, then the Italian for to log in, then say \emph{registrarsi} again
\end{itemize}

\begin{tcolorbox}[breakable,colback=teal!4,colframe=teal!35!black,title={Before you move on}]
What does registrarsi mean? (\textquotedblleft{}To sign up, to register\textquotedblright{}.) Say it once more.
\end{tcolorbox}
